\unnumchapter{KẾT LUẬN, KIẾN NGHỊ}

\textbf{Kết quả của báo cáo định kỳ lần 1.}
Báo cáo đã xác định bài toán trợ lý desktop hỗ trợ quản lý Google Calendar bằng tiếng Việt, làm rõ phạm vi nguyên mẫu, trình bày kiến thức nền về xử lý tiếng nói, hiểu lệnh, suy luận thời gian, OAuth và Calendar API, đồng thời đề xuất yêu cầu, quy trình nghiệp vụ và kiến trúc mô-đun ban đầu. Đây là kết quả phân tích và thiết kế; báo cáo chưa khẳng định các mô-đun đã được triển khai hoặc có số liệu thực nghiệm.

\textbf{Giới hạn.}
Các lựa chọn về VAD, wake word, STT, TTS và công nghệ desktop vẫn cần được kiểm chứng trên môi trường mục tiêu. Chưa có tập dữ liệu tiếng Việt có gán nhãn để đánh giá; quy tắc cho câu lệnh mơ hồ, mặc định thời lượng và thao tác với sự kiện lặp cũng cần được thống nhất. Kết quả nhận dạng và độ trễ thực tế chưa được đo trong giai đoạn này.

\textbf{Kiến nghị và hướng tiếp theo.}
Giai đoạn kế tiếp nên ưu tiên nguyên mẫu microphone--VAD--STT và bộ ca kiểm thử tiếng Việt; sau đó hiện thực phân tích ý định, chuẩn hóa thời gian và OAuth trên lịch thử nghiệm riêng. Chỉ thêm chức năng ghi/sửa/xóa sau khi đã có cơ chế xác nhận và kiểm thử lỗi. Các chỉ tiêu hiệu năng cần được chốt cùng giảng viên hướng dẫn sau khi đo đường cơ sở, đồng thời rà soát bảng phân công và thông tin bìa theo tình hình thực tế.
